% =============================================================================
\section{Methodology}
\label{sec:method}
% =============================================================================

\subsection{Preliminaries: Setwise Ranking}

Given a query $q$ and a set of $n$ candidate documents $\mathcal{D} = \{d_1, d_2, \ldots, d_n\}$ from a first-stage retriever, the goal is to produce a ranking $\pi$ that orders documents by decreasing relevance to $q$. We focus on the top-$k$ ranking problem, where only the top $k$ positions need to be correctly determined ($k \ll n$).

The setwise approach~\citep{zhuang2024setwise} defines a \emph{comparison function} $f_{\text{best}}$ that, given $q$ and a subset $S \subseteq \mathcal{D}$ of size $c+1$ (where $c$ is the \texttt{num\_child} parameter), returns the most relevant document:
\begin{equation}
f_{\text{best}}(q, S) = \arg\max_{d \in S} \text{relevance}(q, d)
\label{eq:fbest}
\end{equation}

This comparison function serves as the comparator in a sorting algorithm (heapsort or bubblesort) to produce the top-$k$ ranking. Two scoring modes are supported:
\begin{itemize}[leftmargin=*]
\item \textbf{Generation}: The LLM generates the passage label (e.g., ``A''); the decoded output determines the winner.
\item \textbf{Likelihood}: The logits of candidate passage label tokens are compared directly. For T5 we use direct next-token logits over single-character labels; for Qwen3/Qwen3.5 we score short teacher-forced continuations such as \texttt{Passage A}, since causal tokenization does not guarantee that label letters are stable single tokens. Both variants avoid autoregressive decoding.
\end{itemize}

\textbf{Prompt template (standard):}
\begin{quote}
\small\texttt{Given a query "\{query\}", which of the following passages is the most relevant one to the query?}\\
\texttt{Passage A: "\{doc\_1\}"}\\
\texttt{Passage B: "\{doc\_2\}"}\\
$\cdots$\\
\texttt{Output only the passage label of the most relevant passage:}
\end{quote}


\subsection{Bottom-Up Setwise Ranking}
\label{sec:bottomup}

We define a reverse comparison function $f_{\text{worst}}$ that selects the \emph{least relevant} document:
\begin{equation}
f_{\text{worst}}(q, S) = \arg\min_{d \in S} \text{relevance}(q, d)
\label{eq:fworst}
\end{equation}

\textbf{Prompt template (reverse):}
\begin{quote}
\small\texttt{Given a query "\{query\}", which of the following passages is the \textbf{least} relevant one to the query?}\\
\texttt{Passage A: "\{doc\_1\}"}\\
\texttt{Passage B: "\{doc\_2\}"}\\
$\cdots$\\
\texttt{Output only the passage label of the \textbf{least} relevant passage:}
\end{quote}

The prompt change is minimal: only two words differ (``least'' vs. ``most''). This ensures that any observed differences in accuracy or bias are attributable to the cognitive task of identifying the worst versus best document, rather than confounding prompt effects.

\textbf{Algorithm}: We adapt both heapsort and bubblesort for bottom-up selection:
\begin{itemize}[leftmargin=*]
\item \emph{Bottom-up heapsort}: We build a \emph{min-heap} where $f_{\text{worst}}$ serves as the comparator, placing the least relevant document at the root. We then run a complete heapsort with the worst-selection comparator: each extraction pulls the worst remaining document to the back of the array, and the survivors at the front are guaranteed to be the top-$k$ in their final positions. Crucially, every comparison made during both heap construction and extraction uses the worst-selection prompt.
\item \emph{Bottom-up bubblesort}: Instead of bubbling the best document to the top, we sink the worst document to the bottom via forward passes. We run $n-1$ passes (rather than $n-k$) so that the top-$k$ region is itself sorted by the same worst-selection signal, rather than left in heap order with a stray best-selection clean-up at the end.
\end{itemize}

\textbf{Efficiency}: Bottom-up heapsort requires $\sim$$n-1$ min-extractions (each with $O(\log_c n)$ heapify cost) using only the worst-selection comparator. For the typical IR setting ($k=10$, $n=100$), this is $\sim$273 comparisons in our experiments versus $\sim$77 for \topdown{}-Heap---a $\sim$$3.5\times$ increase. Bottom-up bubblesort requires $\sim$1683 comparisons in our setting, more than an order of magnitude over \topdown{}-Heap. Consequently, \bottomup{} is more costly for top-$k$ when $k \ll n$, and is primarily valuable for: (i) effectiveness analysis of worst-selection versus best-selection, (ii) as an ensemble component with potentially complementary error patterns, and (iii) scenarios where $k$ is close to $n$ (full ranking).

\textbf{Hypothesis}: We hypothesize that LLMs may exhibit higher accuracy when identifying clearly \emph{irrelevant} documents compared to distinguishing among highly relevant ones. This is because irrelevance detection (``this document about cooking is not relevant to a quantum physics query'') often involves categorical topic mismatch that is easy to detect, whereas best-selection among topically relevant candidates requires nuanced assessment of relevance degree. As we will see in \S\ref{sec:rq1}, this hypothesis is decisively refuted on our 9 models.


\subsection{Dual-End Setwise Ranking}
\label{sec:dualend}

We define a \emph{dual comparison function} $f_{\text{dual}}$ that simultaneously returns both the most and least relevant documents from a candidate set:
\begin{equation}
\resizebox{0.90\columnwidth}{!}{$
f_{\text{dual}}(q, S) = \bigl(\arg\max_{d \in S} \text{rel}(q,d),\ \arg\min_{d \in S} \text{rel}(q,d)\bigr)
$}
\label{eq:fdual}
\end{equation}

\textbf{Prompt template (dual-end):}
\begin{quote}
\small\texttt{Given a query "\{query\}", which of the following passages is the most relevant and which is the least relevant to the query?}\\
\texttt{Passage A: "\{doc\_1\}"}\\
\texttt{Passage B: "\{doc\_2\}"}\\
$\cdots$\\
\texttt{Output the most relevant passage label and the least relevant passage label, separated by a comma:}
\end{quote}

\textbf{Key insight}: When selecting the best document from a set, the LLM internally compares all candidates. It implicitly identifies not only the winner but also the loser---this information is simply never requested. By asking for both, we extract a second ranking decision from the same cognitive process with minimal overhead.

\textbf{Token overhead analysis}: The dual-end prompt adds approximately 15 tokens to the instruction (``\emph{and which is the least relevant to the query}'' and the output format specification). The output grows by approximately 1--2 tokens (the second label). For a typical setwise comparison with 4 passages of 128 tokens each, the total prompt is $\sim$600 tokens. The 15-token overhead is $\sim$2.5\%, while the information gain is $\sim$$74$--$90\%$ depending on window size (see Table~\ref{tab:infobits}).

\textbf{Scoring modes}: Dual-end supports two scoring approaches:
\begin{itemize}[leftmargin=*]
\item \textbf{Generation mode}: The LLM generates text containing both labels (e.g., ``Best: A, Worst: C''). A robust cascading parser handles multiple output formats including ``Passage A, Passage C'', bracketed labels (\texttt{[A]}), numeric labels, and verbose responses. This is the primary mode for Qwen3/Qwen3.5.
\item \textbf{Likelihood mode}: For T5 we read the full softmax over passage label tokens from a single forward pass and take the maximum as best and minimum as worst. For Qwen3/Qwen3.5 we use the same heuristic but score short continuations like \texttt{Passage A} with teacher forcing. We are explicit that this dual-end likelihood path is a best-only-distribution proxy rather than an exact joint $\Pr(\textnormal{Best}{=}X,\textnormal{Worst}{=}Y\mid q,S)$ model. We use it as an efficiency probe and do not over-claim its information-theoretic content.
\end{itemize}

\textbf{Output parsing}: Reliable extraction of both labels from the generated text is critical. We implement a cascading parser:
(1) First, try the explicit ``Best: X, Worst: Y'' format via regex.
(2) If that fails, look for ``Passage X'' patterns.
(3) As a fallback, extract standalone single characters (A--W) from the output.
(4) If parsing fails entirely, default to the first and last passage labels and log a warning.
The success rate of the primary parser on different models is reported in \S\ref{sec:rq2}.


\subsection{Bidirectional Ensemble}
\label{sec:bidir}

We run both \topdown{} and \bottomup{} ranking independently on the same candidate set, producing two rankings $\pi_{\text{td}}$ and $\pi_{\text{bu}}$. These are then fused into a final ranking using rank aggregation.

\textbf{Motivation}: At design time we hypothesized that the error patterns of top-down and bottom-up rankings would be complementary. Top-down errors might occur when the LLM incorrectly identifies a mediocre document as the best in some comparison (false positive at the top), while bottom-up errors might occur when the LLM fails to identify a truly irrelevant document as the worst (false negative at the bottom). Since these error types could be different, fusion might correct errors that either individual method makes. The empirical result (\S\ref{sec:rq3}) is that this complementarity does not materialize because BottomUp's recency-driven errors are too systematic to cancel out, not because the principle of fusion fails.

We evaluate three fusion methods:
\begin{itemize}[leftmargin=*]
\item \textbf{Reciprocal Rank Fusion (RRF)}~\citep{cormack2009reciprocal}: $\text{score}(d) = \sum_{i \in \{\text{td}, \text{bu}\}} \frac{1}{k_\text{rrf} + \text{rank}_i(d)}$ with $k_\text{rrf}=60$. RRF is widely used for combining heterogeneous ranking signals and is robust to score scale differences.
\item \textbf{CombSUM}: Normalize ranks to $[0,1]$ scores per run ($s_i(d) = (n - \text{rank}_i(d) + 1) / n$), then sum. This gives equal weight to both methods.
\item \textbf{Weighted}: $\text{score}(d) = \alpha \cdot s_{\text{td}}(d) + (1-\alpha) \cdot s_{\text{bu}}(d)$ where $\alpha$ controls the relative weight. We experiment with $\alpha \in \{0.3, 0.5, 0.7, 0.9\}$.
\end{itemize}

\textbf{Efficiency}: \bidir{} requires $2\times$ the LLM calls of standard setwise (one full pass each for top-down and bottom-up). This is comparable to permutation voting with $p=2$~\citep{zhuang2024setwise}, which also doubles the cost but provides order-based rather than direction-based complementary information. A key experimental question is which form of complementarity (direction vs. order) yields greater improvement per unit of additional cost.
